\section{Step 6 Probe after an exact fill}\label{png-load-s06}
Use the 115-byte worked file with capacity 115. Once the main buffer fills, there is no room for a further byte there. Implement \texttt{load\_probe} by requesting one byte into the supplied private stack slot at offset 64. A nonzero result takes the failure-closing path: a positive result means extra data, and a negative result means an I/O failure. Zero continues to the EOF check.

The selected capacity bounds the file, not merely the part the program chooses to keep. Without this probe, a valid 115-byte file and a longer file sharing that prefix would look the same to the loader. The probe destination must not be the byte immediately after the accepted buffer.

Use \texttt{s06-1} and stop before \texttt{load\_probe\_ready}.

\paragraph{Observe} The exact-fit file produces a one-byte request with result zero. The accepted 115 bytes and the following guard stay unchanged.

\paragraph{Check} What does a probe result of one establish?
\begin{itemize}
  \item[A.] The file was exactly one byte long.
  \item[B.] The current PNG has one image row.
  \item[C.] The file exceeds the selected capacity.
\end{itemize}
